\clearpage
\section{Wstęp}
\label{chap:wstep}

\subsection{Motywacja i znaczenie problemu}
\label{sec:motywacja}

Dane opisujące transakcje, pacjentów, klientów lub wyniki pomiarów często mają postać tabeli, w~której każdej obserwacji odpowiada wiersz, a~poszczególnym właściwościom --- kolumny. W~zadaniu klasyfikacji na podstawie takiego opisu należy przewidzieć etykietę należącą do skończonego zbioru klas. Przykładami są wykrywanie nadużyć, ocena ryzyka kredytowego i~przewidywanie ponownej hospitalizacji. Wspólna organizacja danych nie oznacza jednak jednorodności problemów: kolumny mogą mieć różne typy, skale i~znaczenie, a~obserwacje mogą być ze sobą powiązane~\cite{borisov2022survey}.

Zastosowanie sieci neuronowej do danych tabelarycznych wymaga uzasadnienia zarówno jakością predykcji, jak i~kosztem obliczeń. Punktem odniesienia są metody liniowe, klasyfikatory sąsiedztwa oraz zespoły drzew. Na wynik porównania wpływają także przygotowanie danych, ustawienia uczenia, dobór hiperparametrów i~procedura walidacji. Dlatego ocena architektury wymaga odpowiednio dobranych modeli bazowych oraz spójnych warunków eksperymentu~\cite{gorishniy2021revisiting}.

W~pracy skupiono się na trzech trudnościach: małej liczbie przykładów uczących, dużej liczbie cech w~stosunku do liczby przykładów oraz niezrównoważeniu klas. W~ostatnim przypadku wysoki odsetek poprawnych predykcji może współwystępować z~niską wykrywalnością zdarzeń rzadkich. Z~kolei przy małej próbie i~wielu cechach szczególnego znaczenia nabiera ocena nadmiernego dopasowania. Porównanie metod zostanie zatem uzupełnione analizą ważenia klas i~regularyzacji.

\subsection{Cel i zakres pracy}
\label{sec:cel-zakres}

Celem pracy jest zbadanie przydatności znanych z~literatury metod głębokiego uczenia do klasyfikacji danych tabelarycznych. Ocena będzie oparta na porównaniu z~wybranymi algorytmami konwencjonalnymi oraz na badaniach mechanizmów uwrażliwiania na rzadkie klasy i~ograniczania nadmiernego dopasowania. Praca ma charakter badawczy; nie zakłada konieczności skonstruowania nowej architektury ani wykazania przewagi określonej rodziny metod.

Przyjęto siedem metod neuronowych: MLP, TabNet, FT-Transformer, RealMLP, ModernNCA, TabM i~TabPFN-3.5. Cztery odniesienia konwencjonalne stanowią regresja logistyczna, KNN, XGBoost i~CatBoost. Wybrane metody różnią się architekturą, reprezentacją danych i~sposobem uczenia. RealMLP pozwala uwzględnić znaczenie ustawień całej procedury uczenia, ModernNCA --- uczonej reprezentacji sąsiedztwa, a~TabPFN-3.5, jako przyklad modelu fundamentalnego --- pretrenowania na innych zadaniach.

Badanie obejmie dwanaście zbiorów: Adult, Covertype, Credit Card Fraud, HIGGS w~wybranej wersji podzbioru, Amazon Employee Access, Santander Customer Satisfaction, Online Shoppers Purchasing Intention, Diabetes 130-US Hospitals, Helena, Statlog German Credit, Ionosphere i~Bioresponse. Zachowano zadania klasyfikacji binarnej i~wieloklasowej. Zbiory dobrano celowo, aby uwzględnić zróżnicowane liczności, typy cech i~rozkłady klas; nie stanowią one losowej próby wszystkich zastosowań uczenia maszynowego.

Porównanie główne będzie uzupełnione dwoma badaniami kontrolowanymi. Pierwsze dotyczy ważenia klas przy niezmienionych pozostałych hiperparametrach. Drugie obejmuje usuwanie wybranych składników regularyzacji MLP przy różnych liczbach przykładów treningowych. Aktywne uczenie pozostaje kierunkiem możliwego rozszerzenia, a~nie częścią wymaganą do udzielenia odpowiedzi na przyjęte pytania.

\subsection{Pytania badawcze}
\label{sec:pytania-badawcze}

Przyjęty cel uszczegółowiono w~pięciu pytaniach. Każde odnosi się do zadeklarowanych danych i~warunków oceny, a~nie do wszystkich możliwych modeli i~zastosowań.

\begin{enumerate}
    \item \textbf{PB1:} Jaką jakość predykcji osiągają wybrane metody neuronowe i~konwencjonalne w~dwunastu zadaniach klasyfikacji? Podstawą odpowiedzi będą wyniki poszczególnych zadań i~różnice między procedurami ocenianymi na wspólnych populacjach testowych.
    \item \textbf{PB2:} Jak strojenie hiperparametrów i~łączenie konfiguracji zmieniają jakość oraz koszt? Porównanie obejmie konfigurację domyślną, konfigurację wybraną na walidacji oraz zespół konfiguracji. Warunki te pozostaną rozdzielone w~raportowaniu.
    \item \textbf{PB3:} Jak ważenie klas wpływa na rozpoznawanie zadeklarowanej rzadkiej klasy? Oceniane będą jakość uporządkowania przykładów oraz precyzja i~czułość przy progach wyznaczonych bez dostępu do etykiet testowych.
    \item \textbf{PB4:} Jak dropout i~zanik wag wpływają na generalizację MLP przy różnej liczbie przykładów treningowych? Odpowiedź wymaga porównań z~ustaloną recepturą odniesienia, wspólnych podzbiorów i~kontroli pozostałych ustawień.
    \item \textbf{PB5:} Jaki kompromis jakości, czasu i~pamięci uzyskują badane procedury na zadeklarowanym sprzęcie? Koszt znalezienia rozwiązania zostanie oddzielony od kosztu wykorzystania wybranych modeli i~wykonania predykcji.
\end{enumerate}

Ocenie będą podlegały zarówno poprawa, jak i~pogorszenie jakości oraz przypadki, w~których metoda nie ukończy obliczeń. Brak statystycznego potwierdzenia różnicy nie będzie interpretowany jako równoważność metod.

\subsection{Wkład własny i granice wnioskowania}
\label{sec:wklad}

Wkład własny obejmuje syntezę literatury, uzasadnienie doboru metod i~zadań, projekt porównania oraz integrację odtwarzalnego środowiska eksperymentalnego. Zaimplementowano główną ścieżkę uruchamiania metod, zapisu predykcji i~odtwarzania raportów. Z~uwagi na badawczy charakter pracy, integracja w miarę możliwości korzysta z~istniejących bibliotek i~implementacji; nie oznacza autorstwa oryginalnych algorytmów. Jej poprawność należy rozpatrywać oddzielnie od empirycznej skuteczności metod.

Planowany wkład empiryczny obejmuje porównanie jakości i~kosztu oraz oba badania mechanizmów. Na etapie przygotowania tego tekstu końcowa kampania nie została wykonana. Dostępne sprawdzenia infrastruktury nie są przedstawiane jako wyniki naukowego porównania. Wnioski dotyczące skuteczności zostaną sformułowane dopiero na podstawie zidentyfikowanych artefaktów eksperymentów.

Porównanie kompletnych rozwiązań nie pozwala samo w~sobie przypisać różnicy jakości pojedynczej warstwie, reprezentacji lub regularyzatorowi. Taki wniosek wymaga odpowiednio dopasowanej interwencji. Podobnie jeden podział zewnętrzny nie opisuje zmienności między wszystkimi możliwymi podziałami, a~foldy walidacji wewnętrznej nie zwiększają liczby niezależnych zbiorów w~porównaniu między zadaniami~\cite{demsar2006comparisons}.

\subsection{Organizacja pracy}
\label{sec:organizacja}

Rozdział~\ref{chap:stan-badan} przedstawia właściwości danych tabelarycznych, trudne reżimy klasyfikacji oraz stan badań porównawczych. Wyjaśnia, jakie informacje wynikają z~benchmarków, a~jakie wymagają badań kontrolowanych. W~rozdziale~\ref{chap:metody} opisano mechanizmy wybranych metod i~uzasadniono ich role w~porównaniu.

Rozdział~\ref{chap:metodyka} opisuje podziały danych, strojenie, metryki i~projekt badań kontrolowanych. Rozdział~\ref{chap:srodowisko} przedstawia implementację środowiska eksperymentalnego. Rozdziały~\ref{chap:wyniki} i~\ref{chap:mechanizmy} przygotowano jako szkielety prezentacji wyników; zostaną uzupełnione po wykonaniu eksperymentów. Planowana dyskusja końcowa odniesie wyniki do literatury i~ograniczeń badania.